\section{Conclusion and Future Work}

We presented the composable multi-level intermediate representation and transformations that underpin tensor code generation in \MLIR.
This work heavily leverages \MLIR's progressivity design principle, which it helped to shape by designing and implementing multiple instances of progressive lowering.
The approach is characterized by a \emph{transformation-oriented IR design}: doing away with legality analyses and applicability checks on low-level IR, relying systematically on the progressive decomposition of carefully designed abstractions.
The resulting design is modular and built with optionality in mind; abstractions span data structures and control flow with both functional (SSA form) and imperative (side-effecting) semantics; they serve as generic building blocks for retargetable tensor compilers.
Transformations are systematically applied as compositions of declarative patterns.
This allows implementing advanced forms of pass fusion, known to alleviate the dreaded phase-ordering problem.
Early single threaded CPU code generation shows strong performance, even in the absence of systematic tuning.

Multiple extensions and generalizations of this work are actively being pursued:
\begin{itemize}
    \item Functionally inspired abstractions for parallel semantics on immutable tensors will extend our results to multi-threaded CPUs and beyond.
    \item Generalization of the approach to other hardware targets thanks to our retargetable vector abstractions.
    \item More powerful vector abstractions including generalized masking and warp-level GPU programming.
    \item More advanced data structures and iterators than dense or sparse to enable cross-pollination with other fields such as databases.
    \item Autotuning infrastructure to systematically find transformation compositions and parameters that produce code running at a high fraction of peak, for any hardware.
    \item Controllable framework for selecting and composing transformation patterns, making production compilers effectively extensible with domain-specific code generation strategies.
\end{itemize}

We believe these contributions will further demystify the field of high-performance compilation,
give a high-level of control to users of the system and provide a strong basis on which future research and near-peak performance production systems will be built.
